\section{Nonescape and ordinary ball bounds}\label{sec:geometry}

We apply the weighted count in two ways. Cubes about zero control
the cusp; small boxes away from the coordinate hyperplanes control
ordinary neighborhoods in the lattice space.

\begin{proposition}\label{prop:tightness}
Every sequence of the packet probabilities with $Q\to\infty$
is tight. More precisely, for $0<\theta<1$,
\begin{equation}\label{eq:cusp}
 \limsup_{Q\to\infty}\mu_*
 \{\Lambda:\exists\,0\ne v\in\Lambda,\ \|v\|_\infty<\theta\}
       \le C_n\theta^n.
\end{equation}
Every subsequential probability limit is $A_n$-invariant.
\end{proposition}

\begin{proof}
For $f=\one_{[-\theta,\theta]^n}$, the first $n-1$ logarithms
in \eqref{eq:product-fiber} are at most $\log\theta$.
The last coordinate bounds the sum of their gaps from
$\log\theta$ by $\log(\theta^n/z)$. Thus
\[
 \cV_f(z)=\frac{2^n}{(n-1)!}
           \bigl(\max\{0,\log(\theta^n/z)\}\bigr)^{n-1}.
\]
Put $X=\theta^nQ$ and $U=Q^{3/4}$. Up to an $n$-dependent
constant, Proposition~\ref{prop:unfolding} reduces the first
moment to
\[
 \frac1{\sqrt D\,\kappaK}
 \sum_{\N\mathfrak a\le X}
     P(\mathfrak a)\left(\log\frac X{\N\mathfrak a}\right)^{n-1}.
\]
If $X<1$, the sum is empty. Otherwise its logarithmic factor is
at most $(\log Q)^{n-1}$, since $\theta<1$. The terms with
$\N\mathfrak a\le\min(U,X)$ contribute at most
$C_nQ^{-1/4}(\log Q)^{n-1}$, by applying
Proposition~\ref{prop:weighted-count} at $x=y=U$.
When $X>U$, define the truncated cumulative sum
\[
 A_U(t)=\sum_{U<\N\mathfrak a\le t}P(\mathfrak a),
 \qquad U\le t\le X.
\]
Writing each logarithmic power as an integral and applying Tonelli gives
\begin{align*}
 \sum_{U<\N\mathfrak a\le X}
   P(\mathfrak a)\left(\log\frac X{\N\mathfrak a}\right)^{n-1}
 &= (n-1)\int_U^X A_U(t)
      \left(\log\frac Xt\right)^{n-2}\frac{dt}{t} \\
 &\ll_n \frac{\kappaK}{q}
      \int_U^X\left(\log\frac Xt\right)^{n-2}\dd t
 \ll_n \frac{\kappaK}{q}X.
\end{align*}
Here the same proposition at $x=y=t$ bounds $A_U(t)$, and the
last inequality follows from
$\int_0^1(\log(1/v))^{n-2}\dd v<\infty$.
There is no lower boundary term, since $A_U(U)=0$, and the
logarithmic weight vanishes at $X$. If $X\le U$, this part is absent.
Dividing by $\sqrt D\,\kappaK$ gives $O_n(\theta^n)$.
Since $E_f\ge1$ whenever the displayed cusp event occurs,
this proves \eqref{eq:cusp}.

Mahler's compactness criterion \cite[Section~3, Theorem~2]{Mahler1946} says that the
complement of this cusp event is compact. Given a desired
mass error, choose $\theta$ first and then $Q$ sufficiently large.
The finitely many earlier packet supports are compact by
Lemma~\ref{lem:packet-model}, so the whole sequence is tight.
The lattice space is locally compact and second countable;
weak compactness for tight probabilities gives probability
subsequential limits. Invariance passes to them by testing
against continuous compactly supported functions.
\end{proof}

For $r>0$, let
\[
 B(r)=\{h\in\SL_n(\R):
             \max_{i,j}|h_{ij}-\delta_{ij}|<r\}.
\]

\begin{proposition}\label{prop:ball-bound}
Let $\mu$ be a weak limit of packet probabilities along
$Q\to\infty$. For every compact $\Omega\subset X_n$ there are
$C_\Omega,r_\Omega>0$ such that
\begin{equation}\label{eq:ball-bound}
 \mu(xB(r))\le C_\Omega r^n
       \qquad(x\in\Omega,\ 0<r<r_\Omega).
\end{equation}
\end{proposition}

\begin{proof}
First fix a compact $V\subset(\R^\times)^n$ and $b_0\ge1$.
For $w\in V$ and small $r$, take
$f=\one_{w+[-b_0r,b_0r]^n}$. Each coordinate has a fixed sign,
is uniformly separated from zero, and has logarithmic width
$O_{V,b_0}(r)$. Its absolute product lies within $Cr$ of
$z_w=\prod_j|w_j|$, so
\[
 \cV_f(z)\le Cr^{n-1}
       \one_{[z_w-Cr,z_w+Cr]}(z)
\]
for some $C\ge1$ depending only on $V,b_0,n$.
Choose $r_0$ small enough that $z_w>4Cr_0$ on $V$.
Use Proposition~\ref{prop:weighted-count} with
\[
 x=(z_w+2Cr)Q,\qquad y=4CrQ.
\]
This slightly enlarged interval contains every contributing
ideal norm. Uniformly in $w$, for sufficiently large $Q$ it has
$Q^{3/4}\le x\le A_0Q$ with a fixed $A_0\ge1$, and
$\eta x\le y\le x$ with $\eta=4Cr/A_0\le1$, after decreasing
$r_0$ if necessary. For every fixed $r\in(0,r_0)$,
Proposition~\ref{prop:unfolding} therefore gives, eventually,
\begin{equation}\label{eq:vector-box}
 \int E_{\one_{w+[-b_0r,b_0r]^n}}\dd\mu_*
       \ll_{V,b_0,n}
       \frac{r^{n-1}}{\sqrt D\,\kappaK}
       \frac{\kappaK rQ}{q}=r^n .
\end{equation}
Here $Q=q\sqrt D$. The implicit constant is independent of $r$. The threshold for $Q$
may depend on $r$, which is all that will be needed.

We next choose a suitable lattice vector at each $x\in\Omega$.
Every full lattice has a vector with no zero coordinate:
for any basis, its combination with coefficients
$(1,N,\ldots,N^{n-1})$ has coordinates that are nonzero
polynomials in $N$, so an integer $N$ can avoid their finitely
many roots. Near each center choose a continuous local lift to basis
matrices, and continue the chosen vector using its fixed integer
basis coefficients. On a sufficiently small neighborhood with compact
closure in that chart, its coordinates remain bounded and bounded away
from zero. A finite cover of $\Omega$ by such neighborhoods therefore
gives choices $w_x$ lying in one compact
$V\subset(\R^\times)^n$. This pointwise selection need not be
continuous or measurable.

For $h$ in $B(r)$, each coordinate of $w_xh-w_x$ has absolute
value at most $r\sum_i|(w_x)_i|$. Thus one fixed $b_0$ makes
every lattice in $xB(r)$ contain a nonzero vector in
$w_x+[-b_0r,b_0r]^n$. Equation~\eqref{eq:vector-box} bounds
the packet mass of $xB(r)$ by $C_\Omega r^n$.
This set is open, so Portmanteau gives
\[
 \mu(xB(r))\le\liminf_i\mu_i(xB(r))\le C_\Omega r^n.
\]
The argument holds for every fixed $x$ and $r$ with the same
constant on $\Omega$. It proves \eqref{eq:ball-bound} for
all $0<r<r_\Omega$, without any exchange of shrinking-radius
and packet limits.
\end{proof}
